\documentclass[11pt]{article}

\usepackage[a4paper, top=0.5cm, left=2cm, right=2cm,bottom=2cm]{geometry}
\usepackage{amsfonts}
\usepackage{amsmath}
\usepackage{mathtools}
\usepackage{amssymb}
\usepackage{multicol}
\usepackage{fullpage}
\usepackage{graphicx}
\usepackage{float}
\usepackage{geometry}
\usepackage{polynom}
\geometry{top=1.2cm}
\title{\small Facultad de Matemática, Astronomía, Física y Computación}
\author{Simulacro Segundo Parcial\\ GURI - La Bisagra, conduccion del CEIMAF}
\date{29 de Agosto de 2026}
\begin{document}
\maketitle
\begin{enumerate}
    \item 
    \begin{enumerate}
        \item Calcular las siguientes expresiones:
        \begin{enumerate}
        \item Trabajemos la forma de reescribir la división de fracciones
        \begin{gather*}
            \frac{\left(\frac{2}{3}\right)^{-2}}{8^{\frac{2}{3}}} \xrightarrow{\text{Reescribamoslo}} \left(\frac{2}{3}\right)^{-2} \cdot \frac{1}{8^{\frac{2}{3}}}
        \end{gather*}
        ahora planteamos las propiedades de la potencia sobre las fracciones cuando la potencia es negativa y cuando es fraccionaria:
        \begin{gather*}
            \left(\frac{2}{3}\right)^{-2}\cdot \frac{1}{8^{\frac{2}{3}}}= \left(\frac{3}{2}\right)^2 \cdot \frac{1}{\left(\sqrt[3]{8}\right)^2} 
        \end{gather*}
        Resolvemos primero la raiz en el denominador y ahora las potencias distribuyendo el exponente:
        \begin{gather*}
            \implies \frac{3^2}{2^2}\cdot \frac{1}{2^2} = \frac{9}{4} \cdot \frac{1}{4} = \frac{9 \cdot 1}{4 \cdot 4} = \frac{9}{16}
        \end{gather*}
        Finalmente la fracción como es irreducible es la solución de la operación.
        \item Identifiquemos los números complejos que están actuando en esta operación, sea $z_1=1+2i$ y $z_2=i$, ahora resolvamos por partes antes de realizar el producto:
        \begin{gather*}
            \overline{z_1}=\overline{1+2i}=1-2i
        \end{gather*}
        \begin{gather*}
            z_2^3=i^3=i^2 \cdot i = -1 \cdot i = -i
        \end{gather*}
        Ahora operemos todo junto:
        \begin{gather*}
            \overline{z_1} \cdot z_2^3 = (1-2i) \cdot (-i)
        \end{gather*}
        Apliquemos la propiedad distributiva y resolvamos
        \begin{gather*}
            1 \cdot (-i) + (-2i) \cdot (-i)= -i + 2i^2\\
            -i + 2 (-1)= -i-2
        \end{gather*}
        \item Comencemos trabajando con los denominadores dentro de la raiz, siempre separando en términos y no olvidando que para sumar y restar fracciones se deben igualar los denominadores:
        \begin{gather*}
            \sqrt{\frac{2}{5}+\frac{1}{1+\frac{1}{1+\frac{1}{2}}}}=\sqrt{\frac{2}{5}+\frac{1}{1+\frac{1}{\frac{2}{2}+\frac{1}{2}}}}=\sqrt{\frac{2}{5}+\frac{1}{1+\frac{1}{\frac{3}{2}}}}
        \end{gather*}
        Recordemos que al tener una fracción en el denominador se invierte y continuamos operando:
        \begin{gather*}
            \sqrt{\frac{2}{5}+\frac{1}{1+\frac{2}{3}}}=\sqrt{\frac{2}{5}+\frac{1}{\frac{3}{3}+\frac{2}{3}}}=\sqrt{\frac{2}{5}+\frac{1}{\frac{5}{3}}}=\sqrt{\frac{2}{5}+\frac{3}{5}}=\sqrt{\frac{5}{5}}
        \end{gather*}
        Terminamos de resolver la raiz, no olvidando que en este caso la raiz es positiva, entonces el resultado es positivo 
        \begin{gather*}
            \implies \sqrt{1}=1
        \end{gather*}
        \end{enumerate}
        \item  Completar la tabla indicando con una cruz a que conjunto/s pertenece cada uno de los resultados obtenidos en las operaciones del ítem anterior.
            \begin{table}[H]
                \centering
                \resizebox{0.98\textwidth}{!}{
                \begin{tabular}{|c|c|c|c|c|c|c|}
                \hline
                \ & $\mathbb{N}$(Naturales) & $\mathbb{Z}$(Enteros) & $\mathbb{Q}$(Racionales)& $\mathbb{I}$(Irracionales)& $\mathbb{R}$(Reales) & $\mathbb{C}$(Complejos) \\ \hline
                I: 9/16 & \ & \ & X & \ & X & \ \\ \hline
                II: $-2-i$ & \ & \ & \ & \ & \ & X \\ \hline
                III: $1$ & X & X & X & \ & X & \ \\ \hline
            \end{tabular}
                }
            \end{table}
    \end{enumerate}
    \item Para resolver operaciones con polinomios lo ideal es trabajar por separado con algunos términos y luego terminar de resolverlos juntos
    \begin{enumerate}
        \item si separamos por términos, mi prioridad es resolver primero el producto:
        \begin{itemize}
            \item $Q(x)\cdot S(x)$ Desarrollemos la propiedad distributiva dentro de los polinomios:
            \begin{align*}
                Q(x)\cdot S(x) &= (x^2-2x) \cdot (x+1)=x^2\cdot x + x^2\cdot 1 -2x\cdot x-2x\cdot 1\\
                &=x^{2+1}+x^2-2x^{1+1}-2x= x^3+x^2-2x^2-2x\\
                &=x^3-x^2-2x
            \end{align*}
            Llamemos $M(x)$ a este polinomio obtenido por el producto de Q y S
            \item $x \cdot M(x)$ Nuevamente distribuimos en el polinomio obtenido
            \begin{align*}
                x\cdot M(x) &=x\cdot(x^3-x^2-2x)\\
                &=x\cdot x^3+x\cdot(-x^2)+x\cdot(-2x)\\
                &=x^4-x^3-2x^2
            \end{align*}
            Llamemos a este polinomio $U(x)$.
            \item Por último nos queda operar con ambos términos:
            \begin{align*}
                P(x)-U(x) &=(x^3-10-4x)-(x^4-x^3-2x^2)\\
                &= x^3-10-4x - x^4+x^3+2x^2
            \end{align*}
            Agrupamos y seguimos operando
            \begin{align*}
                =-x^4+x^3+x^3+2x^2-4x-10\\
                =-x^4+2x^3+2x^2-4x-10
            \end{align*}
            Que es el resultado final!
        \end{itemize}
        \item Resolvemos la división como se conoce teniendo en cuenta los espacios, es decir, que operemos en la resta con los coeficientes que acompañan a las variables de igual grado
        \begin{center}
            \polylongdiv[style=D]{x^3-4x-10}{x^2-2x}
        \end{center}
        Como podemos ver:
        \begin{itemize}
            \item \textbf{Cociente:} $x+2$
            \item \textbf{Resto:} $-10$
        \end{itemize}
        \item Sea $T(x)=Kx^{583}+9x^{264}+2x-7$ y como $S(x)$ es un binomio de grado uno con coeficiente principal 1, podemos aplicar el teorema del resto:\\\\
        Conocemos que $T(x)$ por ser polinomio se puede reescribir de la forma:
        \begin{align*}
            T(x)=Kx^{583}+9x^{264}+2x-7=& S(x) \cdot Q(x) + R(x)\\
            &=(x+1)\cdot Q(x)+5
        \end{align*}
        Restamos 5 a ambos términos y ahora sabemos que el resto tiene que ser 0
        \begin{align*}
            \implies Kx^{583}+9x^{264}+2x-7-5=& S(x) \cdot Q(x) + 5-5\\
            &\implies Kx^{583}+9x^{264}+2x-12=& S(x) \cdot Q(x) \\
            &\implies Kx^{583}+9x^{264}+2x-12=& (x+1) \cdot Q(x) 
        \end{align*}
        Ahora por el teorema del resto, evaluamos $x$ en $(-1)$ y trabajamos con la ecuación:
        \begin{align*}
            \implies K(-1)^{583}+9(-1)^{264}+2(-1)-12= ((-1)+1) \cdot Q(x)
        \end{align*}
        No olvidemos que cuando el exponente es par el resultado obtenido va a ser siempre positivo y cuando es impar, si la base es negativa entonces el resultado va a ser negativo, además que la potencia de 1 siempre va a ser 1, en este caso hay que tener cuidado con el signo
        \begin{align*}
            \implies & K(-1)+9\cdot 1+-2-12= 0 \cdot Q(x)\\
            &-K+9-2+12=0\\
            &19=K
        \end{align*}
        Y así obtenemos el valor de $K$ que buscábamos
    \end{enumerate}
    \item Definición de variables y planteo del sistema. Para modelar matemáticamente la situación, comenzamos asignando variables a las incógnitas que nos pide el problema:
        \begin{itemize}
            \item Sea $m$ la cantidad de entradas vendidas a menores de 10 años (número de menores asistentes).
            \item Sea $a$ la cantidad de entradas vendidas a adultos (número de adultos asistentes).
        \end{itemize}
    A partir de los datos del enunciado, podemos establecer dos relaciones:
    \begin{enumerate}
        \item \textbf{Capacidad total:} Se vendieron todas las entradas para una capacidad de $160$ personas, por lo que la suma de asistentes debe ser:
        \begin{equation}
            m + a = 160
        \end{equation}
        \item \textbf{Recaudación total:} Cada menor abonó \$500 y cada adulto \$3000, acumulando un total de \$230000:
        \begin{equation}
            500m + 3000a = 230000
        \end{equation}
    \end{enumerate}
    Por lo tanto, el sistema de ecuaciones que describe el problema es:
    \begin{equation*}
        \begin{cases}
            m + a = 160 \\
            500m + 3000a = 230000
        \end{cases}
    \end{equation*}
    Pasamos a trabajar con la resolución del sistema. Podemos simplificar la ecuación de recaudación dividiendo ambos miembros por $500$:
    \begin{align*}
        \frac{500m + 3000a}{500} &= \frac{230000}{500} \\
        m + 6a &= 460
    \end{align*}
    Ahora el sistema equivalente resulta:
    \begin{equation*}
        \begin{cases}
            m + a = 160 \quad &\text{(1)} \\
            m + 6a = 460 \quad &\text{(2)}
        \end{cases}
    \end{equation*}
    Despejamos $m$ de la ecuación (1):
    \begin{equation*}
        m = 160 - a
    \end{equation*}
    Sustituimos la expresión de $m$ en la ecuación (2):
    \begin{align*}
        (160 - a) + 6a &= 460 \\
        160 + 5a &= 460 \\
        5a &= 460 - 160 \\
        5a &= 300 \\
        a &= \frac{300}{5} \\
        a &= 60
    \end{align*}
    Con el valor de $a = 60$, calculamos la cantidad de menores $m$:
    \begin{equation*}
        m = 160 - 60 = 100
    \end{equation*}
    Entonces podemos concluir que asistieron \textbf{100 menores de 10 años} y \textbf{60 adultos} al concierto.
    \item Ejercicio 4:
    \begin{enumerate}
        \item  Dominio y valores donde no está definida la ecuación:
        Una expresión algebraica fraccionaria no está definida para aquellos valores de la variable que anulen el denominador, puesto que la división por cero no está permitida en los números reales.
    \begin{equation*}
        \implies x^4 - 5x^2 + 4 = 0
    \end{equation*}
    Una expresión algebraica fraccionaria no está definida para aquellos valores de la variable que anulen el denominador, puesto que la división por cero no está permitida en los números reales.
    \end{enumerate}
\end{enumerate}
\end{document}